\section{Introduction}
\label{sec:intro}

Satisfiability for TCTL asks whether a formula has a timed model at all, rather
than whether a given one satisfies it. A tableau answers the question by
building an AND--OR graph whose nodes carry a set of goals and a zone, and by
looking for a successful marking of that graph~\cite{Clarke99,tableau_ctl}. The
graph is finite because zones are compared modulo a max-constant abstraction,
and its size is bounded by $2^{O(n)}\cdot M^{O(n^{2})}$ for a formula of size
$n$ over constants up to $M$.

That bound has two factors, and the folklore reading is that the second is the
one to attack: refine the abstraction, widen more aggressively, subsume zones
against one another. We began there and got nowhere, so we measured instead.
Counting distinct \emph{discrete} states against total nodes separates the two
factors cleanly, and on the family that costs the most the ratio is $1.00$.
Every node is a distinct combination of outstanding obligations, carrying one
zone. Nothing in the symbolic machinery can help, because the symbolic part is
not doing the multiplying.

What is multiplying is a product of independent obligations. A conjunction of
$k$ bounded eventualities over disjoint atoms has $2^{k}$ combinations of which
are still outstanding, and the search enumerates all of them, though the $k$
questions never interact. The same shape is what a verification engineer writes
when checking several properties of a system at once.

\paragraph{Contribution.} We show the enumeration is avoidable, and by how
much.

\begin{itemize}
\item A \emph{decomposition theorem} (\S\ref{sec:theory}): if a node's goals,
      context and pursuits split into groups mentioning disjoint symbols, the
      node is realisable exactly when each group is. Two directions, one by
      projecting a model and one by taking a product of timed automata in which
      delays synchronise and edges interleave. The theorem needs no condition
      relating the groups' clocks in the zone, which we had expected it to and
      state explicitly because the expectation is natural and wrong.
\item An algorithm (\S\ref{sec:algorithm}) that applies it at every node, with
      a cache keyed on the component, so a component recurring under many
      parents is decided once. Two side conditions carry the weight: clock
      ownership must be read off the formulas rather than the live pursuits,
      and the symbols that couple two groups include data variables, not only
      atomic propositions. Getting either wrong makes an unsatisfiable formula
      satisfiable, and we report both as measured failures.
\item A measurement-driven account (\S\ref{sec:problem}, \S\ref{sec:eval}) of
      what does \emph{not} work: conflict-driven search, zone subsumption,
      coarser $LU$ abstraction, and four more. Each was implemented and
      measured rather than argued about, and each is reported with the
      structural reason it failed.
\end{itemize}

\paragraph{Results.} On conjunctions of independent obligations the growth
changes from $2^{k}$ to $O(k)$: at $k=20$, $4\,195\,583$ nodes and $31.4$
seconds become $240$ nodes and under ten milliseconds. On a suite of $5000$
formulas across ten families the decomposing search returns exactly the same
verdicts as the baseline, and it costs between two and six per cent on formulas
that do not decompose. Independently of the decomposition, profiling removed a
quartic step and two redundant closures from the zone layer, which is worth a
factor of seven on its own.

\paragraph{A bug, and how it was found.} The evaluation harness was built to
carry expected verdicts wherever a window-containment argument settles them.
It reported $39$ wrong answers on a $700$-case sweep. The cause was in the
calculus rather than the implementation: the delay rule split a clock only at
its own window endpoints, so the onsets whose demand is unsatisfiable sat inside
a cell that also held satisfiable ones. We give the fix in
\S\ref{sec:engineering}. It is the reason we think a satisfiability tool of this
kind needs an oracle-carrying benchmark and not only a regression suite.
